% !TEX root = main.tex
\section*{1. 현재 Introduction은 무엇을 설명하는가?}
이 부분은 현재 1.1~1.4의 내용을 순서대로 풀어 정리한 것입니다.
추가할 내용이나 평가 의견은 다음 부분과 구분합니다.

\subsection*{1.1 제품정보의 신뢰성과 공급망 근거}
\textbf{제품의 환경정보를 설명하려면 앞선 공급자의 자료와 그 자료를 사용한 계산을 함께 확인해야 한다는 내용입니다.}

원문은 DPP가 제품의 특성·환경정보·이력을 제공하는 체계라는 설명으로 시작합니다.
EU 배터리 규정을 예로 들어, 제조사가 탄소발자국과 재활용 원료 함량 등의 정보를 제공해야 한다고 설명합니다.

이 정보를 만들려면 제조사는 셀·모듈 공급자 등 앞선 공급망 참여자로부터 자료를 받아야 합니다.
원료를 여러 제품에 나누어 쓰거나 다른 원료와 함께 가공하면, 받은 자료를 각 제품의 정보에 어떻게 반영했는지도 확인해야 합니다.

공급자의 증빙이 올바르게 발급됐다는 사실만으로, 그 자료를 사용해 계산한 제품정보까지 맞다고 할 수는 없습니다.
원문은 \textbf{증빙의 발급을 확인하는 일과 생산 과정의 계산·배분을 확인하는 일을 구분}합니다.

동시에 기업은 상세 수치와 어떤 원료를 받아 어떤 제품을 만들었는지를 모두 공개하기 어렵다고 설명합니다.
따라서 거래가 정해진 규칙을 따랐는지 확인하는 권한과, 상세 이력을 열람하는 권한을 나누려 합니다.

\textbf{다음 절로 이어지는 질문:} 이처럼 거래를 검증했더라도, 나중에 입력 자료 자체에서 오류가 발견되면 어떻게 해야 할까요?

\subsection*{1.2 정상 검증을 보완하는 사후 대응}
\textbf{거래 당시 계산이 맞았더라도, 그 계산에 사용한 자료에서 나중에 오류가 발견될 수 있다는 내용입니다.}

원료에 관한 증빙이 믿을 만한지는 외부의 공급·검사·감사 절차에서 판단합니다.
원장의 거래 검증은 그 자료를 사용한 계산·배분·이전이 규칙을 따랐는지 확인합니다.

따라서 잘못된 자료를 넣었어도 계산 과정 자체는 규칙에 맞을 수 있습니다.
생산과 이전이 진행된 뒤 증빙의 오류가 발견되면, 그 자료를 앞으로도 사용할 수 있게 둘지 다시 판단해야 합니다.

사용 중단이 결정됐을 때 최초 기록에만 문제가 있다고 표시해서는 충분하지 않습니다.
이미 다른 기업에서 그 기록을 사용해 새 기록을 만들었다면, 이후 거래나 Claim 발급에는 그 새 기록이 쓰이기 때문입니다.

원문은 비공개 거래 연결을 조사해 관련 기록을 찾고, 그 기록의 다음 사용을 원장에서 거부하도록 하자고 제안합니다.
\textbf{무엇이 문제인지 판단하고 사용 중단을 결정하는 주체는 감사자이며, 시스템은 기록을 찾고 차단하는 절차를 지원합니다.}

이 절의 마지막 질문은 정보 공개 없이 최종 제품의 주장과 공급망 자료를 연결하고, 나중의 중단 결정을 후속 기록에도 적용하려면 어떤 기록과 권한이 필요한가입니다.

\subsection*{1.3 선행연구와 본 연구의 과제}
\textbf{검증·추적·동결을 각각 처음 제안하는 것이 아니라, 이를 같은 공급망 이력에 연결하려는 연구라고 설명합니다.}

원문은 네 가지 관련 접근을 소개합니다.
Babel 등은 생산 단계의 자격증명과 원료 지분의 이전·소비를 다룹니다.
Udokwu와 Craß는 확인된 입력으로 Sustainability Claim을 증명하는 체계를 다룹니다.

Sahai 등은 분할·합류·가공 과정의 문서 관계와 추적을 다루며 입력 문서 참조를 원장에 남깁니다.
CAP는 정책에 따른 정보 열람과 동결, 실제 거래와 연결된 감사 자료를 지원합니다.

그다음 원문은 최종 제품의 Claim을 확인하는 일과, 문제 출처의 후속 사용을 조사하는 일을 연결합니다.
전자는 Claim이 어떤 회계 상태와 조건에 근거했는지 확인하고, 후자는 그 근거가 어떤 기록으로 이어졌는지 찾습니다.

두 절차가 서로 다른 정보를 가리켜서는 안 됩니다.
\textbf{거래에서 검증한 입력·출력, Claim 발급에 사용한 상태, 감사자가 복원한 연결, 원장이 사용을 막는 대상이 일치해야 한다}는 것이 이 절의 핵심입니다.

\subsection*{1.4 zkDPP의 설계와 기여}
\textbf{앞에서 제시한 검증·추적·차단을 zkDPP에서 어떻게 수행할지 설명합니다.}

먼저 생산·분할·가공·기업 간 이전을 원장에 기록되는 상태 변화로 표현합니다.
참여자는 실제 입력과 출력이 회계·소유권·권한 규칙을 만족함을 증명하고, 원장은 증명과 현재 사용 가능 상태를 확인합니다.

Claim 발급에서는 회계 상태가 발급 정책을 만족하는지 확인한 뒤 해당 기록을 소비합니다.
제품정보와 정책에 연결된 DPP Claim을 등록하므로, 제3자는 상세 거래자료를 열람하지 않고도 제품정보와 발급 이력의 연결을 확인할 수 있습니다.

각 거래는 실제 부모와 이후 소비를 알아내는 데 필요한 정보를 암호화해 남깁니다.
거래 증명은 이 정보가 실제 입력·출력과 맞는지도 검증하므로, 감사자는 그 정보를 열어 검증된 연결을 따라갈 수 있습니다.

원문은 감사 자료를 여는 데이터 키를 별도의 임계 암호화로 보호한다고 설명합니다.
정해진 수 이상의 위원 협력과 외부 승인 절차를 거쳐 감사자에게 키를 전달하며, 초기 키 사본도 올바르게 관리된다고 가정합니다.

사용 중단이 승인되면 감사자는 확정된 원장 시점에서 후속 기록을 찾습니다.
아직 소비되지 않은 기록 전체를 대상으로 삼으며, 완료되지 않은 기업 간 이전도 포함합니다.

동결이 먼저 확정되면 그 기록을 사용한 가공·이전·Claim 발급을 거부합니다.
반대로 Claim 발급이 먼저 확정되었다면 이미 발급된 이력으로 남기며, 동결로 그 발급을 되돌리지는 않습니다.

마지막으로 처음 정한 대상 전체를 같은 확정 시점에서 다시 확인합니다.
모두 미소비·차단 상태이고 처리 결과도 명확해야 전체 완료로 보고하며, 먼저 소비된 대상이나 결과가 불명확한 대상은 사유와 함께 남깁니다.

기여는 두 가지로 정리합니다.
하나는 비공개 거래 검증·Claim 발급·승인된 사후 대응을 연결하는 시스템 설계입니다.
다른 하나는 명시한 가정 아래 정확성·건전성·Privacy와 실패 조건을 분석하는 것입니다.
